Classical navigation stacks remain attractive because mapping, global planning,
local collision avoidance, and safety checks are explicit and independently
testable.  Their local planners can nevertheless freeze, oscillate, or enter
reciprocal deadlock when moving pedestrians block a narrow passage.  End-to-end
learning can encode interaction patterns, but replacing the nominal controller
also changes behavior in routine states where the classical stack is already
competent.

This work asks whether learning can instead be confined to the short interval in
which a classical planner is at risk of failure.  Planning-Guided Recovery and
Rejoin (PGRR) operates as a failure-triggered layer and leaves Nav2 DWB in
control during normal PointGoal navigation.  An observable detector and a
hysteretic state machine invoke recovery only for collision risk, freeze, oscillation, or
deadlock.  The policy never commands wheel velocity: it selects one of 21 local
temporary goals or WAIT, BACKUP, REPLAN, and CONTINUE.  Nav2 executes the
decision, after which PGRR restores the original task goal.  A planning-derived
action mask restricts selectable actions, and an independent braking supervisor
constrains executed motion.

Training exploits simulator state without leaking it to deployment.  A
short-horizon expert rolls out every legal recovery action using the map and
privileged pedestrian state.  Behavior cloning initializes a compact policy,
and DAgger relabels states visited by the learned controller.  At test time, PGRR
receives only LiDAR history, the local path and goal, robot and planner commands,
progress history, planner status, and rule-detector scores.

The contributions are:
\begin{itemize}
  \item a failure-triggered ROS2 hierarchy that preserves classical planning
  authority outside explicitly detected recovery intervals;
  \item a fixed, interpretable recovery action space with planning and sensor
  masks enforced during expert labeling and online policy execution;
  \item an automatically labeled BC--DAgger pipeline that aggregates
  learner-induced recovery states, with controlled ablations of aggregation and
  action masking; and
  \item a predeclared paired five-method protocol over eight
  pedestrian-interaction families, three densities, five held-out repetitions
  per cell, immutable raw streams, and complete-condition statistical analysis.
\end{itemize}
The evaluation reports contact, task completion, timeout, planner failure, and
intervention behavior separately.  Its numerical claims are generated from the
same five-method result table used for confidence intervals and paired tests;
the fail-closed paper build accepts no incomplete or exploratory result stream.
